%! The Statistical Cycle and Types of Data — Unit 1, Lesson 1.1 — Warm-Up (Answer Key)
% ONE PAGE, blank and key. Three items at 12pt (the stats-model size; replaces the earlier
% \large-at-10pt trick), \workrowsep 50pt, item spacing 22pt, so each item gets real
% handwriting room and still fits. Do not add a fourth item — the page is full.
% GRADUAL RELEASE: the warm-up activates prior knowledge before the guided notes. Items 1 and 2
% are the council's own bike-rack survey — the same 60-of-900 numbers the notes table opens
% with — so the survey students meet in row 1 is already half-computed. Item 1's 45% is the
% bus row of the frequency table in row 3 (problem 9), already done for them. Item 2 is the
% scale-up problem 11 and Guided Practice problem 17 both ask for. Item 3 classifies two
% variables WITHOUT deciding them from a question: row 4 — the crux — takes that away and
% makes the QUESTION decide the type before any data exist.
\documentclass[12pt]{article}
\usepackage{saar-article}
\usepackage{saar-key}   % pulls in -boxes; do NOT also load -boxes
\newcommand{\TallMath}[1]{$\displaystyle #1\rule[-1.4em]{0pt}{3.2em}$}
% Word bank strip for every fill-in cluster (user request 2026-09-06). Defined per document;
% the key inherits it and prints the same strip, so heights match.
\newcommand{\wordbank}[1]{\par\vspace{2pt}\noindent\fcolorbox{goldacc}{goldbg}{\parbox{\dimexpr\linewidth-2\fboxsep-2\fboxrule\relax}{\small\textbf{Word bank:} #1}}\par\vspace{4pt}}
% Handwriting room inside every \begin{work} block. Moves the blank and the
% key together, so the two can never drift apart.
\setlength{\workrowsep}{50pt}

\begin{document}

\pageheader{Unit 1, Lesson 1.1}{Warm-Up --- Answer Key}
% NAMESTRIP: no \namedateperiod here — mirrors the blank. NO teachernote here —
% teacher prose lives in the lesson plan.

\vspace{0.15in}

\begin{notesbox}{Spiral Review --- The Student Council's Survey}
\normalsize
Last spring the Riverbend Student Council asked $60$ students, ``How do you usually get to
school?'' Of the $60$ students asked, $27$ said they ride the bus.

\begin{enumerate}[leftmargin=1.8em, itemsep=22pt, topsep=10pt]

  \item What \textbf{percent} of the students asked ride the bus?

        \begin{work}
          \dfrac{27}{60} &= 0.45 \\
                         &= 45\%
        \end{work}

  \item Riverbend has $900$ students in all. Based on that survey, \textbf{predict}
        how many of the $900$ ride the bus.

        \begin{work}
          0.45 \times 900 &= 405 \text{ students}
        \end{work}

  \item The council recorded two things about each student. Classify each variable.
        \wordbank{categorical \;\textbullet\; quantitative \;\textbullet\; the values are labels, not amounts \;\textbullet\; the values are amounts you can average}

        \vspace{4pt}
        \renewcommand{\arraystretch}{2.6}
        \begin{tabularx}{\linewidth}{@{} p{5.4cm} p{3.2cm} >{\raggedright\arraybackslash}X @{}}
        \toprule
        \textbf{Variable} & \textbf{Type} & \textbf{How I know} \\
        \midrule
        How the student gets to school & \ans{categorical} & \ans{the values are labels, not amounts} \\
        Minutes the trip takes         & \ans{quantitative} & \ans{the values are amounts you can average} \\
        \bottomrule
        \end{tabularx}

\end{enumerate}
\end{notesbox}

\end{document}
